\Chapter{10}

\Verse{1}
{
    \zA{话说金荣因人多势众，又兼贾瑞勒令赔了不是，给秦钟磕了头，宝玉方才不吵 闹了。 大家散了学，金荣自己回到家中，越想越气，说：“秦钟不过是贾蓉的小舅
        子，又不是贾家的子孙，附学读书，也不过和我一样。 因他仗着宝玉和他相好，就 目中无人。 既是这样，就该干些正经事，也没的说； 他素日又和宝玉鬼鬼祟祟的，
        只当人家都是瞎子看不见。 今日他又去勾搭人，偏偏撞在我眼里，就是闹出事来， 我还怕什么不成？”}
}
{
    \eA{We will now resume our story. As the persons against Chin Jung were so many
        and their pressure so great, and as, what was more, Chia Jui urged him to
        make amends, he had to knock his head on the ground before Ch'in Chung.
        Pao-yü then gave up his clamorous remonstrances and the whole crowd
        dispersed from school.}
}
{
}

\Verse{2}
{
    \zA{他母亲胡氏听见他咕咕唧唧的，说：“你又要管什么闲事？ 好容 易我和你姑妈说了，你姑妈又千方百计的和他们西府里琏二奶奶跟前说了，你才得
        了这个念书的地方儿。 若不是仗着人家，咱们家里还有力量请的起先生么？ 况且人 家学里茶饭都是现成的，你这二年在那里念书，家里也省好大的嚼用呢!省出来的，
        你又爱穿件体面衣裳。 再者你不在那里念书，你就认得什么薛大爷了？ 那薛大爷一 年也帮了咱们七八十两银子。}
}
{
    \eA{Chin Jung himself returned home all alone, but the more he pondered on the
        occurrence, the more incensed he felt. "Ch'in Chung," he argued, "is simply
        Chia Jung's young brother-in-law, and is no son or grandson of the Chia
        family, and he too joins the class and prosecutes his studies on no other
        footing than that of mine; but it's because he relies upon Pao-yü's
        friendship for him that he has no eye for any one.}
}
{
}

\Verse{3}
{
    \zA{你如今要闹出了这个学房，再想找这么个地方儿，我 告诉你说罢，比登天的还难呢!你给我老老实实的玩一会子睡你的觉去，好多着呢！”
        于是金荣忍气吞声，不多一时，也自睡觉去了。 次日仍旧上学去了，不在话下。 且说他姑妈原给了贾家“玉”字辈的嫡派，名唤贾璜，但其族人那里皆能像宁
        荣二府的家势？ 原不用细说。 这贾璜夫妻守着些小小的产业，又时常到宁荣二府里 去请安，又会奉承凤姐儿并尤氏，所以凤姐儿尤氏也时常资助资助他，方能如此度
        日。}
}
{
    \eA{This being the case, he should be somewhat proper in his behaviour, and
        there would be then not a word to say about it! He has besides all along
        been very mystical with Pao-yü, imagining that we are all blind, and have no
        eyes to see what's up! Here he goes again to-day and mixes with people in
        illicit intrigues;}
}
{
}

\Verse{4}
{
    \zA{今日正遇天气晴明，又值家中无事，遂带了一个婆子，坐上车，来家里走走， 瞧瞧嫂子和侄儿。 说起话儿来，金荣的母亲偏提起昨日贾家学房里的事，从头至尾，
        一五一十，都和他小姑子说了。 这璜大奶奶不听则已，听了怒从心上起，说道：“这 秦钟小杂种是贾门的亲戚，难道荣儿不是贾门的亲戚？
        也别太势利了!况且都做的是 什么有脸的事!就是宝玉也不犯向着他到这个田地。 等我到东府里瞧瞧我们珍大奶 奶，再和秦钟的姐姐说说，叫他评评理！”}
}
{
    \eA{and it's all because they happened to obtrude themselves before my very eyes
        that this rumpus has broken out; but of what need I fear?" His mother, née
        Hu, hearing him mutter; "Why meddle again," she explained, "in things that
        don't concern you? I had endless trouble in getting to speak to your
        paternal aunt;}
}
{
}

\Verse{5}
{
    \zA{金荣的母亲听了，急的了不得，忙说道： “这都是我的嘴快，告诉了姑奶奶，求姑奶奶快别去说罢!别管他们谁是谁非，倘 或闹出来，怎么在那里站的住？
        要站不住，家里不但不能请先生，还得他身上添出 许多嚼用来呢！” 璜大奶奶说道：“那里管的那些个？ 等我说了，看是怎么样！” 也不
        容他嫂子劝，一面叫老婆子瞧了车，坐上竟往宁府里来。 到了宁府，进了东角门，下了车，进去见了尤氏，那里还有大气儿？}
}
{
    \eA{and your aunt had, on the other hand, a thousand and one ways and means to
        devise, before she could appeal to lady Secunda, of the Western mansion; and
        then only it was that you got this place to study in. Had we not others to
        depend upon for your studies, would we have in our house the means
        sufficient to engage a teacher? Besides, in other people's school, tea and
        eatables are all ready and found;}
}
{
}

\Verse{6}
{
    \zA{殷殷勤勤 叙过了寒温，说了些闲话儿，方问道：“今日怎么没见蓉大奶奶？” 尤氏说：“他这 些日子不知怎么了，经期有两个多月没有来。
        叫大夫瞧了，又说并不是喜。 那两日 到下半日就懒怠动了，话也懒怠说，神也发涅。 我叫他：‘你且不必拘礼，早晚不 必照例上来，你竟养养儿罢。
        就有亲戚来，还有我呢。 别的长辈怪你，等我替你告 诉。’ 连蓉哥儿我都嘱咐了，我说：‘你不许累？ 他，不许招他生气，叫他静静儿的 养几天就好了。}
}
{
    \eA{and these two years that you've been there for your lessons, we've likewise
        effected at home a great saving in what would otherwise have been necessary
        for your eating and use. Something has been, it's true, economised; but you
        have further a liking for spick and span clothes. Besides, it's only through
        your being there to study, that you've come to know Mr. Hsüeh!}
}
{
}

\Verse{7}
{
    \zA{他要想什么吃，只管到我屋里来取。 倘或他有个好歹，你再要娶这 么一个媳妇儿，这么个模样儿，这么个性格儿，只怕打着灯笼儿也没处找去呢！’
        他这为人行事儿，那个亲戚长辈儿不喜欢他？ 所以我这两日心里很烦。 偏偏儿的早 起他兄弟来瞧他，谁知那小孩子家不知好歹，看见他姐姐身上不好，这些事也不当
        告诉他，就受了万分委曲也不该向着他说。 谁知昨日学房里打架，不知是那里附学 的学生，倒欺负他，里头还有些不干不净的话，都告诉了他姐姐。}
}
{
    \eA{that Mr. Hsüeh, who has even in one year given us so much pecuniary
        assistance as seventy and eighty taels! And now you would go and raise a row
        in this school-room! why, if we were bent upon finding such another place, I
        tell you plainly, and once for all, that we would find it more difficult
        than if we tried to scale the heavens!}
}
{
}

\Verse{8}
{
    \zA{婶子你是知道的： 那媳妇虽则见了人有说有笑的，他可心细，不拘听见什么话儿都要忖量个三日五夜 才算。 这病就是打这‘用心太过’上得的。
        今儿听见有人欺负了他的兄弟，又是恼， 又是气。 恼的是那狐朋狗友，搬弄是非，调三窝四； 气的是为他兄弟不学好，不上 心念书，才弄的学房里吵闹。
        他为这件事，索性连早饭还没吃。 我才到他那边解劝 了他一会子，又嘱咐了他的兄弟几句，我叫他兄弟到那边府里又找宝玉儿去； 我又
        瞧着他吃了半钟儿燕窝汤，我才过来了。}
}
{
    \eA{Now do quietly play for a while, and then go to sleep, and you'll be ever so
        much better for it then." Chin Jung thereupon stifled his anger and held his
        tongue; and, after a short while, he in fact went to sleep of his own
        accord. The next day he again went to school, and no further comment need be
        made about it;}
}
{
}

\Verse{9}
{
    \zA{婶子，你说我心焦不心焦？ 况且目今又没 个好大夫，我想到他病上，我心里如同针扎的一般!你们知道有什么好大夫没有？”
        金氏听了这一番话，把方才在他嫂子家的那一团要向秦氏理论的盛气，早吓的 丢在爪洼国去了。 听见尤氏问他好大夫的话，连忙答道：“我们也没听见人说什么
        好大夫。 如今听起大奶奶这个病来，定不得还是喜呢。 嫂子倒别教人混治，倘若治 错了，可了不得！” 尤氏道：“正是呢。”}
}
{
    \eA{but we will go on to explain that a young lady related to her had at one
        time been given in marriage to a descendant (of the eldest branch) of the
        Chia family, (whose names were written) with the jade radical, Chia Huang by
        name; but how could the whole number of members of the clan equal in
        affluence and power the two mansions of Ning and Jung? This fact goes, as a
        matter of course, without saying.}
}
{
}

\Verse{10}
{
    \zA{说话之间，贾珍从外进来，见了金氏，便 问尤氏道：“这不是璜大奶奶么？” 金氏向前给贾珍请了安，贾珍向尤氏说：“你让 大妹妹吃了饭去。”
        贾珍说着话便向那屋里去了。 金氏此来原要向秦氏说秦钟欺负 他侄儿的事，听见秦氏有病，连提也不敢提了。 况且贾珍尤氏又待的甚好，因转怒
        为喜的，又说了一会子闲话，方家去了。 金氏去后，贾珍方过来坐下，问尤氏道：“今日他来又有什么说的？”}
}
{
    \eA{The Chia Huang couple enjoyed some small income; but they also went, on
        frequent occasions, to the mansions of Ning and Jung to pay their respects;
        and they knew likewise so well how to adulate lady Feng and Mrs. Yu, that
        lady Feng and Mrs. Yu would often grant them that assistance and support
        which afforded them the means of meeting their daily expenses.}
}
{
}

\Verse{11}
{
    \zA{尤氏答 道：“倒没说什么，一进来脸上倒像有些个恼意似的，及至说了半天话儿，又提起 媳妇的病，他倒渐渐的气色平和了。
        你又叫留他吃饭，他听见媳妇这样的病，也不 好意思只管坐着，又说了几句话就去了，倒没有求什么事。 如今且说媳妇这病，你
        那里寻一个好大夫给他瞧瞧要紧，可别耽误了!现今咱们家走的这群大夫，那里要 得？ 一个个都是听着人的口气儿，人怎么说，他也添几句文话儿说一遍；}
}
{
    \eA{It just occurred on this occasion that the weather was clear and fine, and
        that there happened, on the other hand, to be nothing to attend to at home,
        so forthwith taking along with her a matron, (Mrs. Chia Huang) got into a
        carriage and came over to see widow Chin and her nephew.}
}
{
}

\Verse{12}
{
    \zA{可倒殷勤 的很，三四个人，一日轮流着，倒有四五遍来看脉!大家商量着立个方儿，吃了也 不见效。
        倒弄的一日三五次换衣裳、坐下起来的见大夫，其实于病人无益。” 贾珍 道：“可是这孩子也糊涂，何必又脱脱换换的。 倘或又着了凉，更添一层病，还了 得？
        任凭什么好衣裳，又值什么呢，孩子的身体要紧。 就是一天穿一套新的，也不 值什么。}
}
{
    \eA{While engaged in a chat, Chin Jung's mother accidentally broached the
        subject of the affair, which had transpired in the school-room of the Chia
        mansion on the previous day, and she gave, for the benefit of her young
        sister-in-law, a detailed account of the whole occurrence from beginning to
        end. This Mrs. Huang would not have had her temper ruffled had she not come
        to hear what had happened;}
}
{
}

\Verse{13}
{
    \zA{我正要告诉你：方才冯紫英来看我，他见我有些心里烦，问我怎么了，我 告诉他媳妇身子不大爽快，因为不得个好大夫，断不透是喜是病，又不知有妨碍没
        妨碍，所以我心里实在着急。 冯紫英因说他有一个幼时从学的先生，姓张名友士， 学问最渊博，更兼医理极精，且能断人的生死。 今年是上京给他儿子捐官，现在他
        家住着呢。 这样看来，或者媳妇的病该在他手里除灾也未可定。 我已叫人拿我的名 帖去请了。 今日天晚，或未必来，明日想一定来的。}
}
{
    \eA{but having heard about it, anger sprung from the very depths of her heart.
        "This fellow, Ch'in Chung," she exclaimed, "is a relative of the Chia
        family, but is it likely that Jung Erh isn't, in like manner, a relative of
        the Chia family; and when relatives are many, there's no need to put on
        airs! Besides, does his conduct consist, for the most part, of anything that
        would make one get any face?}
}
{
}

\Verse{14}
{
    \zA{且冯紫英又回家亲替我求他， 务必请他来瞧的。 等待张先生来瞧了再说罢。” 尤氏听说，心中甚喜，因说：“后日是太爷的寿日，到底怎么个办法？” 贾珍
        说道：“我方才到了太爷那里去请安，兼请太爷来家受一受一家子的礼。 太爷因说 道：‘我是清净惯了的，我不愿意往你们那是非场中去。}
}
{
    \eA{In fact, Pao-yü himself shouldn't do injury to himself by condescending to
        look at him. But, as things have come to this pass, give me time and I'll go
        to the Eastern mansion and see our lady Chen and then have a chat with Ch'in
        Chung's sister, and ask her to decide who's right and who's wrong!" Chin
        Jung's mother upon hearing these words was terribly distressed.}
}
{
}

\Verse{15}
{
    \zA{你们必定说是我的生日， 要叫我去受些众人的头，你莫如把我从前注的《阴骘文》给我好好的叫人写出来刻
        了，比叫我无故受众人的头还强百倍呢!倘或明日后日这两天一家子要来，你就在 家里好好的款待他们就是了。 也不必给我送什么东西来。 连你后日也不必来。 你要
        心中不安，你今日就给我磕了头去。 倘或后日你又跟许多人来闹我，我必和你不依。’ 如此说了，今日我是再不敢去的了。 且叫赖升来，吩咐他预备两日的筵席。”}
}
{
    \eA{"It's all through my hasty tongue," she observed with vehemence, "that I've
        told you all, sister-in-law: but please, sister, give up at once the idea of
        going over to say anything about it! Don't trouble yourself as to who is in
        the right, and who is in the wrong; for were any unpleasantness to come out
        of it, how could we here stand on our legs?}
}
{
}

\Verse{16}
{
    \zA{尤氏因叫了贾蓉来：“吩咐赖升照例预备两日的筵席，要丰丰富富的。 你再亲 自到西府里请老太太、大太太、二太太和你琏二婶子来逛逛。 你父亲今日又听见一
        个好大夫，已经打发人请去了，想明日必来。 你可将他这些日子的病症细细的告诉 他。” 贾蓉一一答应着出去了。}
}
{
    \eA{and were we not to stand on our legs, not only would we never be able to
        engage a tutor, but the result will be, on the contrary, that for his own
        person will be superadded many an expense for eatables and necessaries."
        "What do I care about how many?" replied Mrs. Huang; "wait till I've spoken
        about it, and we'll see what will be the result."}
}
{
}

\Verse{17}
{
    \zA{正遇着刚才到冯紫英家去请那先生的小子回来了， 因回道：“奴才方才到了冯大爷家，拿了老爷名帖请那先生去，那先生说是：‘方才
        这里大爷也和我说了，但只今日拜了一天的客，才回到家，此时精神实在不能支持， 就是去到府上也不能看脉，须得调息一夜，明日务必到府。’ 他又说：‘医学浅薄，
        本不敢当此重荐，因冯大爷和府上既已如此说了，又不得不去，你先替我回明大人 就是了。 大人的名帖着实不敢当。’ 还叫奴才拿回来了。
        哥儿替奴才回一声儿罢。”}
}
{
    \eA{Nor would she accede to her sister-in-law's entreaties, but bidding, at the
        same time, the matron look after the carriage, she got into it, and came
        over to the Ning Mansion. On her arrival at the Ning Mansion, she entered by
        the eastern side gate, and dismounting from the carriage, she went in to
        call on Mrs. Yu, the spouse of Chia Chen, with whom she had not the courage
        to put on any high airs;}
}
{
}

\Verse{18}
{
    \zA{贾蓉复转身进去，回了贾珍尤氏的话，方出来叫了赖升，吩咐预备两日的筵席的话。 赖升答应，自去照例料理，不在话下。
        且说次日午间，门上人回道：“请的那张先生来了。” 贾珍遂延入大厅坐下。 茶 毕，方开言道：“昨日承冯大爷示知老先生人品学问，又兼深通医学，小弟不胜钦
        敬。” 张先生道：“晚生粗鄙下士，知识浅陋。 昨因冯大爷示知，大人家第谦恭下士， 又承呼唤，不敢违命。 但毫无实学，倍增汗颜。”}
}
{
    \eA{but gently and quietly she made inquiries after her health, and after
        passing some irrelevant remarks, she ascertained: "How is it I don't see
        lady Jung to-day?" "I don't know," replied Mrs. Yu, "what's the matter with
        her these last few days; but she hasn't been herself for two months and
        more; and the doctor who was asked to see her declares that it is nothing
        connected with any happy event.}
}
{
}

\Verse{19}
{
    \zA{贾珍道：“先生不必过谦，就请先 生进去看看儿妇，仰仗高明，以释下怀。” 于是贾蓉同了进去，到了内室，见了秦 氏，向贾蓉说道：“这就是尊夫人了？”
        贾蓉道：“正是。 请先生坐下，让我把贱内 的病症说一说再看脉如何？” 那先生道：“依小弟意下，竟先看脉，再请教病源为 是。
        我初造尊府，本也不知道什么，但我们冯大爷务必叫小弟过来看看，小弟所以 不得不来。 如今看了脉息，看小弟说得是不是，再将这些日子的病势讲一讲，大家
        斟酌一个方儿。}
}
{
    \eA{A couple of days back, she felt, as soon as the afternoon came, both to
        move, and both even to utter a word; while the brightness of her eyes was
        all dimmed; and I told her, 'You needn't stick to etiquette, for there's no
        use for you to come in the forenoon and evening, as required by
        conventionalities; but what you must do is, to look after your own health.}
}
{
}

\Verse{20}
{
    \zA{可用不可用，那时大爷再定夺就是了。” 贾蓉道：“先生实在高明， 如今恨相见之晚。 就请先生看一看脉息可治不可治，得以使家父母放心。” 于是家
        下媳妇们，捧过大迎枕来，一面给秦氏靠着，一面拉着袖口，露出手腕来。 这先生 方伸手按在右手脉上，调息了至数，凝神细诊了半刻工夫。 换过左手，亦复如是。
        诊毕了，说道：“我们外边坐罢。” 贾蓉于是同先生到外边屋里炕上坐了。 一个婆子端了茶来，贾蓉道：“先生请 茶。”}
}
{
    \eA{Should any relative come over, there's also myself to receive them; and
        should any of the senior generation think your absence strange, I'll explain
        things for you, if you'll let me.' "I also advised brother Jung on the
        subject: 'You shouldn't,' I said, 'allow any one to trouble her; nor let her
        be put out of temper, but let her quietly attend to her health, and she'll
        get all right.}
}
{
}

\Verse{21}
{
    \zA{茶毕，问道：“先生看这脉息还治得治不得？” 先生说：“看得尊夫人脉息， 左寸沉数，左关沉伏，右寸细而无力，右关虚而无神。 其左寸沉数者，乃心气虚而
        生火； 左关沉伏者，乃肝家气滞血亏。 右寸细而无力者，乃肺经气分太虚； 右关虚 而无神者，乃脾土被肝木克制。 心气虚而生火者，应现今经期不调，夜间不寐。
        肝 家血亏气滞者，应胁下痛胀，月信过期，心中发热。}
}
{
    \eA{Should she fancy anything to eat, just come over here and fetch it; for, in
        the event of anything happening to her, were you to try and find another
        such a wife to wed, with such a face and such a disposition, why, I fear,
        were you even to seek with a lantern in hand, there would really be no place
        where you could discover her.}
}
{
}

\Verse{22}
{
    \zA{肺经气分太虚者，头目不时眩 晕，寅卯间必然自汗，如坐舟中，脾土被肝木克制者，必定不思饮食，精神倦怠， 四肢酸软。 据我看这脉，当有这些症候才对。
        或以这个的为喜脉，则小弟不敢闻命 矣。” 旁边一个贴身伏侍的婆子道：“何尝不是这样呢!真正先生说得如神，倒不用 我们说了。
        如今我们家里现有好几位太医老爷瞧着呢，都不能说得这样真切。 有的 说道是喜，有的说道是病； 这位说不相干，这位又说怕冬至前后：总没有个真著话 儿。
        求老爷明白指示指示。”}
}
{
    \eA{And with such a temperament and deportment as hers, which of our relatives
        and which of our elders don't love her?' That's why my heart has been very
        distressed these two days! As luck would have it early this morning her
        brother turned up to see her, but who would have fancied him to be such a
        child, and so ignorant of what is proper and not proper to do? He saw well
        enough that his sister was not well;}
}
{
}

\Verse{23}
{
    \zA{那先生说：“大奶奶这个症候，可是众位耽搁了!要在初次行经的时候就用药治 起，只怕此时已全愈了。 如今既是把病耽误到这地位，也是应有此灾。 依我看起来，
        病倒尚有三分治得。 吃了我这药看，若是夜间睡的着觉，那时又添了二分拿手了。 据我看这脉息，大奶奶是个心性高强、聪明不过的人。 但聪明太过，则不如意事常
        有； 不如意事常有，则思虑太过：此病是忧虑伤脾，肝木忒旺，经血所以不能按时 而至。 大奶奶从前行经的日子问一问，断不是常缩，必是常长的。}
}
{
    \eA{and what's more all these matters shouldn't have been recounted to her; for
        even supposing he had received the gravest offences imaginable, it behoved
        him anyhow not to have broached the subject to her! Yesterday, one would
        scarcely believe it, a fight occurred in the school-room, and some pupil or
        other who attends that class, somehow insulted him;}
}
{
}

\Verse{24}
{
    \zA{是不是？” 这婆 子答道：“可不是!从没有缩过，或是长两日三日，以至十日不等，都长过的。” 先 生听道：“是了，这就是病源了。
        从前若能以养心调气之药服之，何至于此!这如今 明显出一个水亏火旺的症候来。 待我用药看。”}
}
{
    \eA{besides, in this business, there were a good many indecent and improper
        utterances, but all these he went and told his sister!}
}
{
}

\Verse{25}
{
    \zA{于是写了方子，递与贾蓉，上写的 是： 益气养荣补脾和肝汤 人参二钱　白术二钱土炒　云苓三钱　熟地四钱　归身二钱　白芍二钱　川芎
        一钱五分　黄芪三钱　香附米二钱　醋柴胡八分　淮山药二钱炒　真阿胶二钱蛤粉 炒　延胡索 钱半酒炒　炙甘草八分　引用建莲子七粒去心　大枣二枚
        贾蓉看了说：“高明的很。 还要请教先生：这病与性命终久有妨无妨？” 先生 笑道：“大爷是最高明的人：人病到这个地位，非一朝一夕的症候了； 吃了这药，
        也要看医缘了。}
}
{
    \eA{Now, sister-in-law, you are well aware that though (our son Jung's) wife
        talks and laughs when she sees people, that she is nevertheless imaginative
        and withal too sensitive, so that no matter what she hears, she's for the
        most part bound to brood over it for three days and five nights, before she
        loses sight of it, and it's from this excessive sensitiveness that this
        complaint of hers arises.}
}
{
}

\Verse{26}
{
    \zA{依小弟看来，今年一冬是不相干的； 总是过了春分，就可望全愈了。” 贾蓉也是个聪明人，也不往下细问了。
        于是贾蓉送了先生去了，方将这药方子并脉案都给贾珍看了，说的话也都回了 贾珍并尤氏了。 尤氏向贾珍道：“从来大夫不像他说的痛快，想必用药不错的。” 贾
        珍笑道：“他原不是那等混饭吃久惯行医的人，因为冯紫英我们相好，他好容易求 了他来的。 既有了这个人，媳妇的病或者就能好了。 他那方子上有人参，就用前日
        买的那一斤好的罢。”}
}
{
    \eA{Today, when she heard that some one had insulted her brother, she felt both
        vexed and angry; vexed that those fox-like, cur-like friends of his had
        moved right and wrong, and intrigued with this one and deluded that one;
        angry that her brother had, by not learning anything profitable, and not
        having his mind set upon study, been the means of bringing about a row at
        school;}
}
{
}

\Verse{27}
{
    \zA{贾蓉听毕了话，方出来叫人抓药去煎给秦氏吃。 不知秦氏服了此药，病势如何，且听下回分解。 (本章完)}
}
{
    \eA{and on account of this affair, she was so upset that she did not even have
        her early meal. I went over a short while back and consoled her for a time,
        and likewise gave her brother a few words of advice; and after having packed
        off that brother of hers to the mansion on the other side, in search of
        Pao-yü, and having stood by and seen her have half a bowl of birds' nests
        soup, I at length came over.}
}
{
}

\Verse{28}
{
    \zA{}
}
{
    \eA{Now, sister-in-law, tell me, is my heart sore or not? Besides, as there's
        nowadays no good doctor, the mere thought of her complaint makes my heart
        feel as if it were actually pricked with needles! But do you and yours,
        perchance, know of any good practitioner?" Mrs. Chin had, while listening to
        these words, been, at an early period, so filled with concern that she cast
        away to distant lands the reckless rage she had been in recently while at
        her sister-in-law's house, when she had determined to go and discuss matters
        over with Mrs. Ch'in.}
}
{
}

\Verse{29}
{
    \zA{}
}
{
    \eA{Upon hearing Mrs. Yu inquire of her about a good doctor, she lost no time in
        saying by way of reply: "Neither have we heard of any one speak of a good
        doctor; but from the account I've just heard of Mrs. Ch'in's illness, it may
        still, there's no saying, be some felicitous ailment; so, sister-in-law,
        don't let any one treat her recklessly, for were she to be treated for the
        wrong thing, the result may be dreadful!" "Quite so!" replied Mrs. Yu. But
        while they were talking, Chia Chen came in from out of doors, and upon
        catching sight of Mrs. Chin; "Isn't this Mrs. Huang?" he inquired of Mrs.
        Yu; whereupon Mrs. Chin came forward and paid her respects to Chia Chen.
        "Invite this lady to have her repast here before she goes," observed Chia
        Chen to Mrs. Yu;}
}
{
}

\Verse{30}
{
    \zA{}
}
{
    \eA{and as he uttered these words he forthwith walked into the room on the off
        side. The object of Mrs. Chin's present visit had originally been to talk to
        Mrs. Ch'in about the insult which her brother had received from the hands of
        Ch'in Chung, but when she heard that Mrs. Ch'in was ill, she did not have
        the courage to even so much as make mention of the object of her errand.
        Besides, as Chia Chen and Mrs. Yu had given her a most cordial reception,
        her resentment was transformed into pleasure, so that after a while spent in
        a further chat about one thing and another, she at length returned to her
        home. It was only after the departure of Mrs. Chin that Chia Chen came over
        and took a seat. "What did she have to say for herself during this visit
        to-day?" he asked of Mrs. Yu.}
}
{
}

\Verse{31}
{
    \zA{}
}
{
    \eA{"She said nothing much," replied Mrs. Yu. "When she first entered the room,
        her face bore somewhat of an angry look, but, after a lengthy chat and as
        soon as mention of our son's wife's illness was made, this angered look
        after all gradually abated. You also asked me to keep her for the repast,
        but, having heard that our son's wife was so ill she could not very well
        stay, so that all she did was to sit down, and after making a few more
        irrelevant remarks, she took her departure. But she had no request to make.
        To return however now to the illness of Jung's wife, it's urgent that you
        should find somewhere a good doctor to diagnose it for her;}
}
{
}

\Verse{32}
{
    \zA{}
}
{
    \eA{and whatever you do, you should lose no time. The whole body of doctors who
        at present go in and out of our household, are they worth having? Each one
        of them listens to what the patient has to say of the ailment, and then,
        adding a string of flowery sentences, out he comes with a long rigmarole;
        but they are exceedingly diligent in paying us visits; and in one day, three
        or four of them are here at least four and five times in rotation! They come
        and feel her pulse, they hold consultation together, and write their
        prescriptions, but, though she has taken their medicines, she has seen no
        improvement; on the contrary, she's compelled to change her clothes three
        and five times each day, and to sit up to see the doctor;}
}
{
}

\Verse{33}
{
    \zA{}
}
{
    \eA{a thing which, in fact, does the patient no good." "This child too is
        somewhat simple," observed Chia Chen; "for what need has she to be taking
        off her clothes, and changing them for others? And were she again to catch a
        chill, she would add something more to her illness; and won't it be
        dreadful! The clothes may be no matter how fine, but what is their worth,
        after all? The health of our child is what is important to look to! and were
        she even to wear out a suit of new clothes a-day, what would that too amount
        to? I was about to tell you that a short while back, Feng Tzu-ying came to
        see me, and, perceiving that I had somewhat of a worried look, he asked me
        what was up;}
}
{
}

\Verse{34}
{
    \zA{}
}
{
    \eA{and I told him that our son's wife was not well at all, that as we couldn't
        get any good doctor, we couldn't determine with any certainty, whether she
        was in an interesting condition, or whether she was suffering from some
        disease; that as we could neither tell whether there was any danger or not,
        my heart was, for this reason, really very much distressed. Feng Tzu-ying
        then explained that he knew a young doctor who had made a study of his
        profession, Chang by surname, and Yu-shih by name, whose learning was
        profound to a degree; who was besides most proficient in the principles of
        medicine, and had the knack of discriminating whether a patient would live
        or die;}
}
{
}

\Verse{35}
{
    \zA{}
}
{
    \eA{that this year he had come to the capital to purchase an official rank for
        his son, and that he was now living with him in his house. In view of these
        circumstances, not knowing but that if, perchance, the case of our
        daughter-in-law were placed in his hands, he couldn't avert the danger, I
        readily despatched a servant, with a card of mine, to invite him to come;
        but the hour to-day being rather late, he probably won't be round, but I
        believe he's sure to be here to-morrow. Besides, Feng-Tzu-ying was also on
        his return home, to personally entreat him on my behalf, so that he's bound,
        when he has asked him, to come and see her. Let's therefore wait till Dr.
        Chang has been here and seen her, when we can talk matters over!"}
}
{
}

\Verse{36}
{
    \zA{}
}
{
    \eA{Mrs. Yu was very much cheered when she heard what was said. "The day after
        to-morrow," she felt obliged to add, "is again our senior's, Mr. Chia
        Ching's birthday, and how are we to celebrate it after all?" "I've just been
        over to our Senior's and paid my respects," replied Chia Chen, "and further
        invited the old gentleman to come home, and receive the congratulations of
        the whole family. "'I'm accustomed,' our Senior explained, 'to peace and
        quiet, and have no wish to go over to that worldly place of yours; for you
        people are certain to have published that it's my birthday, and to entertain
        the design to ask me to go round to receive the bows of the whole lot of
        you.}
}
{
}

\Verse{37}
{
    \zA{}
}
{
    \eA{But won't it be better if you were to give the "Record of Meritorious Acts,"
        which I annotated some time ago, to some one to copy out clean for me, and
        have it printed? Compared with asking me to come, and uselessly receive the
        obeisances of you all, this will be yea even a hundred times more
        profitable! In the event of the whole family wishing to pay me a visit on
        any of the two days, to-morrow or the day after to-morrow, if you were to
        stay at home and entertain them in proper style, that will be all that is
        wanted; nor will there be any need to send me anything! Even you needn't
        come two days from this; and should you not feel contented at heart, well,
        you had better bow your head before me to-day before you go.}
}
{
}

\Verse{38}
{
    \zA{}
}
{
    \eA{But if you do come again the day after to-morrow, with a lot of people to
        disturb me, I shall certainly be angry with you.' After what he said, I will
        not venture to go and see him two days hence; but you had better send for
        Lai Sheng, and bid him get ready a banquet to continue for a couple of
        days." Mrs. Yu, having asked Chia Jung to come round, told him to direct Lai
        Sheng to make the usual necessary preparations for a banquet to last for a
        couple of days, with due regard to a profuse and sumptuous style. "You go
        by-and-by," (she advised him), "in person to the Western Mansion and invite
        dowager lady Chia, mesdames Hsing and Wang, and your sister-in-law Secunda
        lady Lien to come over for a stroll.}
}
{
}

\Verse{39}
{
    \zA{}
}
{
    \eA{Your father has also heard of a good doctor, and having already sent some
        one to ask him round, I think that by to-morrow he's sure to come; and you
        had better tell him, in a minute manner, the serious symptoms of her ailment
        during these few days." Chia Jung having signified his obedience to each of
        her recommendations, and taken his leave, was just in time to meet the youth
        coming back from Feng Tzu-ying's house, whither he had gone a short while
        back to invite the doctor round. "Your slave," he consequently reported,
        "has just been with a card of master's to Mr. Feng's house and asked the
        doctor to come.}
}
{
}

\Verse{40}
{
    \zA{}
}
{
    \eA{'The gentleman here,' replied the doctor, 'has just told me about it; but
        to-day, I've had to call on people the whole day, and I've only this moment
        come home; and I feel now my strength (so worn out), that I couldn't really
        stand any exertion. In fact were I even to get as far as the mansion, I
        shouldn't be in a fit state to diagnose the pulses! I must therefore have a
        night's rest, but, to-morrow for certain, I shall come to the mansion. My
        medical knowledge,' he went on to observe, 'is very shallow, and I don't
        deserve the honour of such eminent recommendation; but as Mr. Feng has
        already thus spoken of me in your mansion, I can't but present myself. It
        will be all right if in anticipation you deliver this message for me to your
        honourable master;}
}
{
}

\Verse{41}
{
    \zA{}
}
{
    \eA{but as for your worthy master's card, I cannot really presume to keep it.'
        It was again at his instance that I've brought it back; but, Sir, please
        mention this result for me (to master)." Chia Jung turned back again, and
        entering the house delivered the message to Chia Chen and Mrs. Yu; whereupon
        he walked out, and, calling Lai Sheng before him, he transmitted to him the
        orders to prepare the banquet for a couple of days. After Lai Sheng had
        listened to the directions, he went off, of course, to get ready the
        customary preparations; but upon these we shall not dilate, but confine
        ourselves to the next day.}
}
{
}

\Verse{42}
{
    \zA{}
}
{
    \eA{At noon, a servant on duty at the gate announced that the Doctor Chang, who
        had been sent for, had come, and Chia Chen conducted him along the Court
        into the large reception Hall, where they sat down; and after they had
        partaken of tea, he broached the subject. "Yesterday," he explained, "the
        estimable Mr. Feng did me the honour to speak to me of your character and
        proficiency, venerable doctor, as well as of your thorough knowledge of
        medicine, and I, your mean brother, was filled with an immeasurable sense of
        admiration!" "Your Junior," remonstrated Dr. Chang, "is a coarse, despicable
        and mean scholar and my knowledge is shallow and vile!}
}
{
}

\Verse{43}
{
    \zA{}
}
{
    \eA{but as worthy Mr. Feng did me the honour yesterday of telling me that your
        family, sir, had condescended to look upon me, a low scholar, and to favour
        me too with an invitation, could I presume not to obey your commands? But as
        I cannot boast of the least particle of real learning, I feel overburdened
        with shame!" "Why need you be so modest?" observed Chia Chen; "Doctor, do
        please walk in at once to see our son's wife, for I look up, with full
        reliance, to your lofty intelligence to dispel my solicitude!" Chia Jung
        forthwith walked in with him. When they reached the inner apartment, and he
        caught sight of Mrs. Ch'in, he turned round and asked Chia Jung, "This is
        your honourable spouse, isn't it?" "Yes, it is," assented Chia Jung;}
}
{
}

\Verse{44}
{
    \zA{}
}
{
    \eA{"but please, Doctor, take a seat, and let me tell you the symptoms of my
        humble wife's ailment, before her pulse be felt. Will this do?" "My mean
        idea is," remarked the Doctor, "that it would, after all, be better that I
        should begin by feeling her pulse, before I ask you to inform me what the
        source of the ailment is. This is the first visit I pay to your honourable
        mansion; besides, I possess no knowledge of anything; but as our worthy Mr.
        Feng would insist upon my coming over to see you, I had in consequence no
        alternative but to come. After I have now made a diagnosis, you can judge
        whether what I say is right or not, before you explain to me the phases of
        the complaint during the last few days, and we can deliberate together upon
        some prescription;}
}
{
}

\Verse{45}
{
    \zA{}
}
{
    \eA{as to the suitableness or unsuitableness of which your honourable father
        will then have to decide, and what is necessary will have been done."
        "Doctor," rejoined Chia Jung, "you are indeed eminently clear sighted; all I
        regret at present is that we have met so late! But please, Doctor, diagnose
        the state of the pulse, so as to find out whether there be hope of a cure or
        not; if a cure can be effected, it will be the means of allaying the
        solicitude of my father and mother." The married women attached to that
        menage forthwith presented a pillow; and as it was being put down for Mrs.
        Ch'in to rest her arm on, they raised the lower part of her sleeve so as to
        leave her wrist exposed. The Doctor thereupon put out his hand and pressed
        it on the pulse of the right hand.}
}
{
}

\Verse{46}
{
    \zA{}
}
{
    \eA{Regulating his breath (to the pulsation) so as to be able to count the
        beatings, he with due care and minuteness felt the action for a considerable
        time, when, substituting the left hand, he again went through the same
        operation. "Let us go and sit outside," he suggested, after he had concluded
        feeling her pulses. Chia Jung readily adjourned, in company with the Doctor,
        to the outer apartment, where they seated themselves on the stove-couch. A
        matron having served tea; "Please take a cup of tea, doctor," Chia Jung
        observed. When tea was over, "Judging," he inquired, "Doctor, from the
        present action of the pulses, is there any remedy or not?"}
}
{
}

\Verse{47}
{
    \zA{}
}
{
    \eA{"The action of the pulse, under the forefinger, on the left hand of your
        honorable spouse," proceeded the Doctor, "is deep and agitated; the left
        hand pulse, under the second finger, is deep and faint. The pulse, under the
        forefinger, of the right hand, is gentle and lacks vitality. The right hand
        pulse, under my second finger, is superficial, and has lost all energy. The
        deep and agitated beating of the forepulse of the left hand arises from the
        febrile state, due to the weak action of the heart. The deep and delicate
        condition of the second part of the pulse of the left wrist, emanates from
        the sluggishness of the liver, and the scarcity of the blood in that organ.}
}
{
}

\Verse{48}
{
    \zA{}
}
{
    \eA{The action of the forefinger pulse, of the right wrist, is faint and lacks
        strength, as the breathing of the lungs is too weak. The second finger pulse
        of the right wrist is superficial and devoid of vigour, as the spleen must
        be affected injuriously by the liver. The weak action of the heart, and its
        febrile state, should be the natural causes which conduce to the present
        irregularity in the catamenia, and insomnia at night; the poverty of blood
        in the liver, and the sluggish condition of that organ must necessarily
        produce pain in the ribs;}
}
{
}

\Verse{49}
{
    \zA{}
}
{
    \eA{while the overdue of the catamenia, the cardiac fever, and debility of the
        respiration of the lungs, should occasion frequent giddiness in the head,
        and swimming of the eyes, the certain recurrence of perspiration between the
        periods of 3 to 5 and 5 to 7, and the sensation of being seated on board
        ship. The obstruction of the spleen by the liver should naturally create
        distaste for liquid or food, debility of the vital energies and prostration
        of the four limbs. From my diagnosis of these pulses, there should exist
        these various symptoms, before (the pulses and the symptoms can be said) to
        harmonise.}
}
{
}

\Verse{50}
{
    \zA{}
}
{
    \eA{But should perchance (any doctor maintain) that this state of the pulses
        imports a felicitous event, your servant will not presume to give an ear to
        such an opinion!" A matron, who was attached as a personal attendant (to
        Mrs. Ch'in,) and who happened to be standing by interposed: "How could it be
        otherwise?" she ventured. "In real truth, Doctor, you speak like a
        supernatural being, and there's verily no need for us to say anything! We
        have now, ready at hand, in our household, a good number of medical
        gentlemen, who are in attendance upon her, but none of these are proficient
        enough to speak in this positive manner. Some there are who say that it's a
        genital complaint; others maintain that it's an organic disease.}
}
{
}

\Verse{51}
{
    \zA{}
}
{
    \eA{This doctor explains that there is no danger: while another, again, holds
        that there's fear of a crisis either before or after the winter solstice;
        but there is, in one word, nothing certain said by them. May it please you,
        sir, now to favour us with your clear directions." "This complaint of your
        lady's," observed the Doctor, "has certainly been neglected by the whole
        number of doctors; for had a treatment with certain medicines been initiated
        at the time of the first occurrence of her habitual sickness, I cannot but
        opine that, by this time, a perfect cure would have been effected. But
        seeing that the organic complaint has now been, through neglect, allowed to
        reach this phase, this calamity was, in truth, inevitable.}
}
{
}

\Verse{52}
{
    \zA{}
}
{
    \eA{My ideas are that this illness stands, as yet, a certain chance of recovery,
        (three chances out of ten); but we will see how she gets on, after she has
        had these medicines of mine. Should they prove productive of sleep at night,
        then there will be added furthermore two more chances in the grip of our
        hands. From my diagnosis, your lady is a person, gifted with a preëminently
        excellent, and intelligent disposition; but an excessive degree of
        intelligence is the cause of frequent contrarieties; and frequent
        contrarieties give origin to an excessive amount of anxious cares. This
        illness arises from the injury done, by worrying and fretting, to the
        spleen, and from the inordinate vigour of the liver; hence it is that the
        relief cannot come at the proper time and season.}
}
{
}

\Verse{53}
{
    \zA{}
}
{
    \eA{Has not your lady, may I ask, heretofore at the period of the catamenia,
        suffered, if indeed not from anaemia, then necessarily from plethora? Am I
        right in assuming this or not?" "To be sure she did," replied the matron;
        "but she has never been subject to anaemia, but to a plethora, varying from
        either two to three days, and extending, with much irregularity, to even ten
        days." "Quite so!" observed the Doctor, after hearing what she had to say,
        "and this is the source of this organic illness! Had it in past days been
        treated with such medicine as could strengthen the heart, and improve the
        respiration, would it have reached this stage? This has now overtly made
        itself manifest in an ailment originating from the paucity of water and the
        vigour of fire;}
}
{
}

\Verse{54}
{
    \zA{}
}
{
    \eA{but let me make use of some medicines, and we'll see how she gets on!" There
        and then he set to work and wrote a prescription, which he handed to Chia
        Jung, the purpose of which was: Decoction for the improvement of
        respiration, the betterment of the blood, and the restoration of the spleen.
        Ginseng, Atractylodes Lancea; Yunnan root; Prepared Ti root; Aralia edulis;
        Peony roots; Levisticum from Sze Ch'uan; Sophora tormentosa; Cyperus
        rotundus, prepared with rice; Gentian, soaked in vinegar; Huai Shan Yao
        root; Real "O" glue; Carydalis Ambigua; and Dried liquorice. Seven Fukien
        lotus seeds, (the cores of which should be extracted,) and two large zizyphi
        to be used as a preparative. "What exalted intelligence!" Chia Jung, after
        perusing it, exclaimed.}
}
{
}

\Verse{55}
{
    \zA{}
}
{
    \eA{"But I would also ask you, Doctor, to be good enough to tell me whether this
        illness will, in the long run, endanger her life or not?" The Doctor smiled.
        "You, sir, who are endowed with most eminent intelligence (are certain to
        know) that when a human illness has reached this phase, it is not a
        derangement of a day or of a single night; but after these medicines have
        been taken, we shall also have to watch the effect of the treatment! My
        humble opinion is that, as far as the winter of this year goes, there is no
        fear; in fact, after the spring equinox, I entertain hopes of a complete
        cure." Chia Jung was likewise a person with all his wits about him, so that
        he did not press any further minute questions.}
}
{
}

\Verse{56}
{
    \zA{}
}
{
    \eA{Chia Jung forthwith escorted the Doctor and saw him off, and taking the
        prescription and the diagnosis, he handed them both to Chia Chen for his
        perusal, and in like manner recounted to Chia Chen and Mrs. Yu all that had
        been said on the subject. "The other doctors have hitherto not expressed any
        opinions as positive as this one has done," observed Mrs. Yu, addressing
        herself to Chia Chen, "so that the medicines to be used are, I think, surely
        the right ones!"}
}
{
}

\Verse{57}
{
    \zA{}
}
{
    \eA{"He really isn't a man," rejoined Chia Chen, "accustomed to give much of his
        time to the practice of medicine, in order to earn rice for his support: and
        it's Feng Tzu-ying, who is so friendly with us, who is mainly to be thanked
        for succeeding, after ever so much trouble, in inducing him to come. But now
        that we have this man, the illness of our son's wife may, there is no
        saying, stand a chance of being cured. But on that prescription of his there
        is ginseng mentioned, so you had better make use of that catty of good
        quality which was bought the other day." Chia Jung listened until the
        conversation came to a close, after which he left the room, and bade a
        servant go and buy the medicines, in order that they should be prepared and
        administered to Mrs. Ch'in.}
}
{
}

\Verse{58}
{
    \zA{}
}
{
    \eA{What was the state of Mrs. Ch'in's illness, after she partook of these
        medicines, we do not know; but, reader, listen to the explanation given in
        the chapter which follows.}
}
{
}
